\chapter{The Foundation of the Teaching}
\label{ch:foundation}

\begin{versebox}
The eighteen root terms ---\\
where are they to be seen?\\
In the Foundation of the Teaching,\\
where every discourse finds its home.
\end{versebox}

\noindent We began this book with eighteen root terms --- the master categories that organize the entire Buddhist canon. Worldly and transcendent. Beings and principles. Knowledge and what is to be known. Seeing and development. The Buddha's own words and the words of others. What should be answered and what should not. Action and its fruit. Wholesome and unwholesome. Permitted and prohibited. And the Netti's own praise of the teaching.\footnote{Pali: \textit{``Tattha aṭṭhārasa mūlapadā kuhiṃ daṭṭhabbā? Sāsanapaṭṭhāne.''} ``Where are the eighteen root terms to be seen? In the Foundation of the Teaching (\textit{sāsanapaṭṭhāna}).'' This final major section of the Nettippakaraṇa grounds the entire hermeneutical system in the actual classification of suttas. The eighteen root terms (\textit{aṭṭhārasa mūlapadā}) in full: worldly / transcendent / both; being-based / principle-based / both; knowledge / knowable / both; seeing / development / both; one's own statement / another's / both; answerable / unanswerable / both; action / result / both; wholesome / unwholesome / both; permitted / prohibited / both; and the praise.}

Now the question is: where do those root terms actually live? Not in the abstract --- but in the suttas themselves. What kinds of sutta exist in the Buddha's word, and how does each kind map to the root terms?

\section*{Sixteen Types of Sutta}

The Netti classifies every discourse in the entire canon into sixteen types. This is not an arbitrary filing system. It grows organically from the three types of defilement (craving, view, and misconduct) and their three corresponding purifications (calm, insight, and good conduct), combined with four stages of spiritual development.\footnote{Pali: \textit{``Saṃkilesabhāgiyaṃ suttaṃ, vāsanābhāgiyaṃ suttaṃ, nibbedhabhāgiyaṃ suttaṃ, asekkhabhāgiyaṃ suttaṃ \ldots taṇhāsaṃkilesabhāgiyaṃ suttaṃ, diṭṭhisaṃkilesabhāgiyaṃ suttaṃ, duccaritasaṃkilesabhāgiyaṃ suttaṃ, taṇhāvodānabhāgiyaṃ suttaṃ, diṭṭhivodānabhāgiyaṃ suttaṃ, duccaritavodānabhāgiyaṃ suttaṃ.''} Sixteen types of sutta, built from four primary types and their combinations. The four primary types: (1) defilement-sharing (\textit{saṃkilesabhāgiya}), (2) habituation-sharing (\textit{vāsanābhāgiya}), (3) penetration-sharing (\textit{nibbedhabhāgiya}), (4) non-trainee-sharing (\textit{asekkhabhāgiya}). These four, combined pairwise and in triples, yield additional types; and the three specific defilement-subtypes plus their three purification-counterparts add six more, for sixteen total.}

The four primary types:

A \textbf{defilement-sharing} sutta describes the dark side --- the mechanisms of suffering, the traps of craving, the consequences of heedlessness. A \textbf{habituation-sharing} sutta describes the bright side of worldly existence --- the rewards of generosity, the fruits of virtue, the pleasures of heaven. A \textbf{penetration-sharing} sutta describes the path itself --- the Noble Eightfold Path, the factors of awakening, the direct knowledge that cuts through ignorance. A \textbf{non-trainee-sharing} sutta describes the state beyond training --- the qualities of an arahant, the peace of final liberation.

These four can be combined. A sutta that describes both the dark side and the bright worldly side is defilement-and-habituation-sharing. One that describes the dark side and the path is defilement-and-penetration-sharing. And so on, through all the logical permutations, until you have sixteen types that cover every conceivable kind of discourse.\footnote{Pali: \textit{``Imāni cattāri suttāni, sādhāraṇāni katāni aṭṭha bhavanti, tāniyeva aṭṭha suttāni sādhāraṇāni katāni soḷasa bhavanti.''} ``These four types of sutta, made into shared [combinations], become eight. Those same eight, made into shared [combinations], become sixteen.'' The three defilement-subtypes (craving, view, misconduct) are also distinguished, along with their purification-counterparts: craving-defilement → calm; view-defilement → insight; misconduct-defilement → good conduct.}

But the Netti does not merely describe these types. It \textit{demonstrates} them by quoting sutta after sutta --- an extraordinary anthology that ranges across the entire Pali canon. The reader who has spent eighteen chapters learning \textit{how} to read the Buddha now gets to \textit{read} the Buddha, with the full apparatus at hand.

\section*{The Dark Anthology: Defilement-Sharing Suttas}

These are the suttas that show the mechanisms of suffering --- vivid, unsettling, and precise:

\begin{versebox}
Blinded by desire, covered by craving's net,\\
shrouded by craving's canopy,\\
bound by the bond of heedlessness ---\\
like fish in the mouth of a trap,\\
they go on to aging and death,\\
as a calf follows its mother.
\end{versebox}

\noindent The image is devastating: we follow aging and death the way a newborn calf follows its mother --- instinctively, helplessly, not knowing where it is going.\footnote{Pali: \textit{``Kāmandhā jālasañchannā, taṇhāchadanachāditā; pamattabandhanā baddhā, macchāva kumināmukhe; jarāmaraṇamanventi, vaccho khīrapakova mātaranti.''} (Ud 64). The first of approximately twenty defilement-sharing suttas quoted in the Sāsanapaṭṭhāna.}

\begin{versebox}
Mind is the forerunner of all states.\\
Mind is chief; all states are mind-made.\\
If with a corrupt mind\\
one speaks or acts ---\\
suffering follows,\\
as the wheel follows the ox's hoof.
\end{versebox}

\noindent The very first verse of the Dhammapada. Everything begins with the mind. Corrupt the source, and everything downstream is poisoned.\footnote{Pali: \textit{``Manopubbaṅgamā dhammā, manoseṭṭhā manomayā; manasā ce paduṭṭhena, bhāsati vā karoti vā; tato naṃ dukkhamanveti, cakkaṃva vahato padanti.''} (Dhp 1).}

\begin{versebox}
When one is sluggish and a great eater,\\
a sleeper who rolls about ---\\
like a great hog fed on grain,\\
that fool enters the womb again and again.
\end{versebox}

\begin{versebox}
As rust arisen from iron\\
eats the very iron from which it arose ---\\
so one's own deeds\\
lead the over-indulgent to a bad destination.
\end{versebox}

\begin{versebox}
Like a thief caught at the breach in the wall,\\
beaten and bound by his own deed ---\\
just so, people in the hereafter,\\
beaten and bound by their own deed.
\end{versebox}

\noindent The images are from everyday life --- a fat pig, rust eating iron, a burglar caught in the act. Nothing abstract. Nothing obscure. The Buddha is pointing at what you can see with your own eyes.\footnote{Pali: ``Sluggish hog'' = Dhp 325 (\textit{middhī yadā hoti mahagghaso ca}). ``Rust from iron'' = Dhp 240 (\textit{ayasāva malaṃ samuṭṭhitaṃ}). ``Thief at the breach'' (\textit{coro yathā sandhimukhe gahīto}). All classified as defilement-sharing suttas.}

\begin{versebox}
Creatures that love happiness ---\\
whoever harms them with a stick,\\
seeking happiness for himself:\\
hereafter, he finds no happiness.
\end{versebox}

\begin{versebox}
When a herd of cattle is crossing a ford\\
and the bull goes crooked,\\
they all go crooked,\\
because the leader has gone crooked.\\[6pt]
Just so among human beings:\\
whoever is held to be the chief ---\\
if he practices unrighteousness,\\
how much more the other people!\\
The whole kingdom suffers\\
when the king is unrighteous.
\end{versebox}

\noindent The political parable is startling --- a teaching about leadership embedded in a teaching about defilement. The corrupt leader doesn't just harm himself; like a bull at a ford, he drags everyone after him.\footnote{Pali: ``Creatures that love happiness'' = Dhp 131 (\textit{sukhakāmāni bhūtāni, yo daṇḍena vihiṃsati}). ``Bull goes crooked'' = AN 4.70 (\textit{gunnaṃ ce taramānānaṃ, jimhaṃ gacchati puṅgavo}).}

\begin{versebox}
The fruit kills the banana tree,\\
the fruit kills the bamboo,\\
the fruit kills the reed ---\\
honor kills the worthless person,\\
as a foal kills the she-mule.
\end{versebox}

\begin{versebox}
An axe is born in a person's mouth,\\
by which a fool cuts himself\\
when he speaks what is ill-spoken.\\[6pt]
Not a sharpened sword,\\
not poison, not deadly venom ---\\
nothing ruins so surely\\
as speech that is ill-spoken.
\end{versebox}

\noindent Your own words are the weapon. The Netti classifies both of these as defilement-sharing suttas --- they show how the mechanisms of suffering work.\footnote{Pali: ``Fruit kills the banana tree'' = AN 4.68 (\textit{phalaṃ ve kadaliṃ hanti}). ``Axe in the mouth'' = SN 1.180 (\textit{purisassa hi jātassa, kuṭhārī jāyate mukhe}).}

\begin{versebox}
Whoever praises the blameworthy,\\
or blames the one who deserves praise ---\\
he gathers misfortune with his mouth;\\
through that misfortune, he finds no happiness.\\[6pt]
This is a small misfortune ---\\
losing one's wealth at dice.\\
Far greater is this misfortune:\\
corrupting one's mind toward the Blessed Ones.
\end{versebox}

\noindent The gambling metaphor cuts deep: you can lose everything at dice and still recover. But corrupt your mind toward the awakened ones? That loss is beyond calculation.\footnote{Pali: \textit{``Yo nindiyaṃ pasaṃsati, taṃ vā nindati yo pasaṃsiyo \ldots ayameva mahantataro kali; yo sugatesu manaṃ padosaye.''} (AN 4.3 / SN 1.180). The Netti quotes the verse about 36 million aeons and 5 additional eons in the Nirabbuda hell --- the consequence of vilifying the noble ones.}

\section*{The Bright Anthology: Habituation-Sharing Suttas}

Now the mirror image. The same Dhammapada verse, but with the polarity reversed:

\begin{versebox}
Mind is the forerunner of all states.\\
Mind is chief; all states are mind-made.\\
If with a serene mind\\
one speaks or acts ---\\
happiness follows,\\
as a shadow that never departs.
\end{versebox}

\noindent One word changed --- corrupt to serene --- and the entire universe shifts.\footnote{Pali: \textit{``Manopubbaṅgamā dhammā, manoseṭṭhā manomayā; manasā ce pasannena, bhāsati vā karoti vā; tato naṃ sukhamanveti, chāyāva anapāyinīti.''} (Dhp 2). A habituation-sharing sutta.}

And then a remarkable prose passage. Mahānāma the Sakyan --- a layman, a householder, living in the crowded city of Kapilavatthu --- comes to the Buddha terrified:

``Lord, Kapilavatthu is prosperous and crowded, the streets are jammed. When I come home at evening, after visiting the Blessed One, I run into elephants, horses, chariots, carts, people. At that moment, Lord, my mindfulness of the Buddha slips. My mindfulness of the Dhamma slips. My mindfulness of the Sangha slips. And I think: `If I were to die right now, what would be my destination? What would be my future course?'\,''

The Buddha's answer is one of the most reassuring passages in the entire canon:

``Do not fear, Mahānāma. Do not fear. Your death will not be a bad one. A noble disciple endowed with four qualities --- verified confidence in the Buddha, the Dhamma, the Sangha, and virtuous conduct --- leans toward nirvana, inclines toward nirvana, slopes toward nirvana. Just as a tree that leans east, when cut at the root, will fall eastward --- just so, Mahānāma, a noble disciple endowed with these four qualities leans toward nirvana.''\footnote{Pali: \textit{``Mahānāmo sakko bhagavantaṃ etadavoca \ldots `Mā bhāyi, mahānāma, mā bhāyi, mahānāma, apāpakaṃ te maraṇaṃ bhavissati \ldots Seyyathāpi, mahānāma, rukkho pācīnaninno pācīnapoṇo pācīnapabbhāro, so mūlacchinno katamena papateyyāti? Yena bhante ninno yena poṇo yena pabbhāroti.\,''} (SN 55.21--22). A habituation-sharing sutta. The direction of a life is set by its habitual inclination, not by its final moment. Mahānāma's anxiety is achingly human; the Buddha's answer is that practice tilts you toward liberation the way gravity tilts a tree.}

The direction of a life is set by its habitual inclination, not by the circumstances of its final moment. You don't have to die in a perfect state of concentration. You just have to have lived leaning the right way.

The habituation-sharing suttas also include extraordinary stories of merit bearing fruit across cosmic time:

\begin{versebox}
Having offered a single flower,\\
for a thousand crores of aeons,\\
among gods and among humans ---\\
by the remainder, he attained final extinguishment.
\end{versebox}

\begin{versebox}
At a green-leaved Bodhi-tree,\\
at a flourishing, full-grown tree,\\
I formed a single perception directed at the Buddha ---\\
mindful, I gained that much.\\[6pt]
Thirty aeons from that day,\\
I do not know a bad destination.\\
The three true knowledges have been realized ---\\
the habituation of that single perception.
\end{versebox}

\noindent A single flower. A single moment of clarity at a tree. And the ripples extend across thirty aeons. That is what habituation means: the small act, repeated or even just sincerely done once, bends the arc of existence.\footnote{Pali: ``Single flower'' = Theragāthā 96 (\textit{ekapupphaṃ cajitvāna, sahassaṃ kappakoṭiyo}). ``Green Bodhi-tree'' = Theragāthā 217 (\textit{assatthe haritobhāse \ldots ekaṃ buddhagataṃ saññaṃ}). Both classified as habituation-sharing suttas. The Netti's anthology continues with stories of stupa-offerings, garland-offerings, and the power of a serene mind directed at the Buddha-field.}

And the bull-ox parable returns, but now in its bright form:

\begin{versebox}
When a herd of cattle is crossing a ford\\
and the bull goes straight,\\
they all go straight,\\
because the leader has gone straight.\\[6pt]
Just so among human beings:\\
whoever is held to be the chief ---\\
if he practices righteousness,\\
how much more the other people!\\
The whole kingdom lives in happiness\\
when the king is righteous.
\end{versebox}

\noindent The same image, reversed. The dark and bright anthologies are mirrors of each other, and the Netti places them side by side so you can see the symmetry.\footnote{Pali: \textit{``Gunnañce taramānānaṃ, ujuṃ gacchati puṅgavo; sabbā tā ujuṃ gacchanti, nette ujuṃ gate sati \ldots sabbaṃ raṭṭhaṃ sukhaṃ seti, rājā ce hoti dhammikoti.''} (AN 4.70, bright version). The Netti quotes both the dark and bright versions, making the point that the same sutta-structure can serve either the defilement or habituation classification depending on its content.}

\section*{The Penetration Anthology: Path Suttas}

These are the most profound --- suttas that describe the actual moment of breaking through:

\begin{versebox}
Freed above, below, everywhere ---\\
not perceiving ``this I am'' ---\\
thus liberated, one has crossed the flood\\
never crossed before, for non-renewed existence.
\end{versebox}

\noindent Four lines. The entire path. ``Not perceiving `this I am'\,'' --- that is the key. Not the destruction of the self, but the cessation of the habit of fabricating one.\footnote{Pali: \textit{``Uddhaṃ adho sabbadhi vippamutto, ayaṃ ahasmīti anānupassī; evaṃ vimutto udatāri oghaṃ, atiṇṇapubbaṃ apunabbhavāyāti.''} (Ud 61). A penetration-sharing sutta.}

And then the extraordinary passage where the Buddha tells Ānanda that the entire path unfolds by natural law, without needing to be willed:

``Ānanda, the virtuous person need not intend: `May freedom from remorse arise in me.' It is the natural law that freedom from remorse arises in the virtuous.''

``The person free from remorse need not intend: `May gladness arise in me.' It is the natural law that gladness arises in one free from remorse.''

``The glad person need not intend: `May rapture arise in me.' It is the natural law that rapture arises in the glad.''

And so on: rapture leads to tranquility, tranquility to happiness, happiness to concentration, concentration to seeing things as they really are, seeing things as they really are to disenchantment, disenchantment to dispassion, dispassion to liberation, liberation to the knowledge-and-vision of liberation. At every step, the same refrain: ``It is the natural law.''\footnote{Pali: \textit{``Sīlavato, ānanda, na cetanā karaṇīyā `kinti me avippaṭisāro jāyeyyā'ti. Dhammatā esā, ānanda, yaṃ sīlavato avippaṭisāro jāyeyya \ldots dhammatā esā ānanda yaṃ vimuttassa vimuttiñāṇadassanaṃ uppajjeyyāti.''} (AN 11.2). One of the most remarkable passages in the entire Pali canon. Each step of the path unfolds by ``natural law'' (\textit{dhammatā}) --- not by effortful intention but by the natural unfolding of conditions. The practitioner sets the initial conditions; the rest unfolds like a river flowing downhill. A penetration-sharing sutta.}

The practitioner sets the initial conditions --- virtue --- and the rest unfolds like a river flowing downhill.

\begin{versebox}
When truths become manifest\\
to the ardent, meditating seeker ---\\
then all his doubts vanish,\\
for he understands each truth with its cause.
\end{versebox}

\begin{versebox}
When truths become manifest\\
to the ardent, meditating seeker ---\\
then all his doubts vanish,\\
for he knows the destruction of conditions.
\end{versebox}

\noindent The same verse, twice, with one word changed: first, understanding the \textit{cause}; then, knowing the \textit{destruction} of conditions. The progression from seeing how things arise to seeing how they cease.\footnote{Pali: \textit{``Yadā have pātubhavanti dhammā, ātāpino jhāyato brāhmaṇassa; athassa kaṅkhā vapayanti sabbā, yato pajānāti sahetudhammanti \ldots yato khayaṃ paccayānaṃ avedīti.''} (Ud 1, first and second verses). Penetration-sharing suttas.}

\begin{versebox}
``What should one cut to sleep in peace?\\
What should one cut to grieve no more?\\
What is the one thing, Gotama,\\
whose killing you approve?''\\[6pt]
``Cut anger, and you sleep in peace.\\
Cut anger, and you grieve no more.\\
Anger --- with its poisoned root\\
and honeyed tip, O brahmin ---\\
that is the killing the noble ones approve.\\
Cut that, and you grieve no more.''
\end{versebox}

\noindent The brahmin's question is surgical: what single thing should be destroyed? The Buddha's answer is equally precise. And notice the phrase ``poisoned root, honeyed tip'' --- anger feels \textit{good} in the moment, which is exactly what makes it so dangerous.\footnote{Pali: \textit{``Kiṃsu chetvā sukhaṃ seti \ldots Kodhaṃ chetvā sukhaṃ seti \ldots kodhassa visamūlassa, madhuraggassa brāhmaṇa; vadhaṃ ariyā pasaṃsanti, taṃ hi chetvā na socatīti.''} (SN 1.187). A penetration-sharing sutta. The phrase \textit{visamūla} (``poisoned root'') combined with \textit{madhur'agga} (``sweet-tipped'') is one of the Buddha's most psychologically acute observations.}

\begin{versebox}
``What should one strike when it arises?\\
What should one dispel when it is born?\\
What should the wise one abandon?\\
What is the realization that brings happiness?''\\[6pt]
``One should strike anger when it arises.\\
One should dispel lust when it is born.\\
The wise one should abandon ignorance.\\
The realization of truth brings happiness.''
\end{versebox}

\noindent Three defilements, three responses, precisely matched: anger must be struck \textit{immediately} (it escalates); lust must be \textit{dispelled} (it seduces); ignorance must be \textit{abandoned} (it runs deep). And the fruit of all three: the realization of truth.\footnote{Pali: \textit{``Kiṃsu hane uppatitaṃ \ldots Kodhaṃ hane uppatitaṃ, rāgaṃ jātaṃ vinodaye; avijjaṃ pajahe dhīro, saccābhisamayo sukhoti.''} A penetration-sharing sutta. The three defilements are not treated identically: each requires its own specific response.}

\begin{versebox}
As if struck by a sword,\\
as if your head were on fire ---\\
a monk should wander mindfully,\\
to abandon sensual lust.\\[6pt]
As if struck by a sword,\\
as if your head were on fire ---\\
a monk should wander mindfully,\\
to abandon identity-view.
\end{versebox}

\noindent The same urgency, twice. First for lust, then for the view of self. The Netti classifies both as penetration-sharing --- they are about the path, about what must be done \textit{now}.\footnote{Pali: \textit{``Sattiyā viya omaṭṭho, ḍayhamānova matthake; kāmarāgappahānāya, sato bhikkhu paribbaje \ldots sakkāyadiṭṭhippahānāya, sato bhikkhu paribbajeti.''} (SN 1.21 / Theragāthā 39). The repetition with variation --- same urgency, different target --- is itself a teaching device.}

And the magnificent ``One Fine Night'' sutta:

\begin{versebox}
Let one not chase after the past\\
or place hopes in the future.\\
What is past has been left behind;\\
what is future has not been reached.\\[6pt]
But whoever sees clearly\\
the present state as it arises ---\\
immovable, unshakable ---\\
let the wise one cultivate that.\\[6pt]
Today the effort must be made!\\
Who knows --- death may come tomorrow.\\
There is no bargaining\\
with Death and his great army.\\[6pt]
The sage calls one who dwells thus,\\
ardent, tireless, day and night,\\
``the one with one fine night.''
\end{versebox}

\noindent Past: gone. Future: not yet here. The present: the only place where practice is possible. And ``one fine night'' --- a single night of genuine presence is worth more than a lifetime of distraction.\footnote{Pali: \textit{``Atītaṃ nānvāgameyya, nappaṭikaṅkhe anāgataṃ \ldots Taṃ ve bhaddekarattoti, santo ācikkhate munīti.''} (MN 131 / SN). The Bhaddekaratta Sutta --- ``One Fine Night'' or ``An Auspicious Night.'' A penetration-sharing sutta. The Buddha's teaching on presence is not abstract mindfulness advice but an urgent call: \textit{today} is the day.}

\begin{versebox}
Where do lust and hatred come from?\\
Discontent and delight and hair-raising dread ---\\
where do they originate?\\
Where do the mind's thoughts arise,\\
like boys releasing a crow?\\[6pt]
Born from craving, self-arisen,\\
like the trunk-born shoots of a banyan tree,\\
clinging to sense-pleasures far and wide\\
like a creeper spread through the forest.\\[6pt]
Those who understand their origin\\
dispel them --- listen, spirit!\\
They cross this flood so hard to cross,\\
never crossed before, for non-renewed existence.
\end{versebox}

\noindent The image of boys releasing a crow --- thoughts let loose like a captured bird, flying in all directions --- is unforgettable.\footnote{Pali: \textit{``Rāgo ca doso ca kutonidānā \ldots kumārakā dhaṅkamivossajanti \ldots Ye naṃ pajānanti yatonidānaṃ, te naṃ vinodenti suṇohi yakkha; te duttaraṃ oghamimaṃ taranti, atiṇṇapubbaṃ apunabbhavāyāti.''} (SN 1.237 / Sn 273). Addressed to a \textit{yakkha} (spirit). A penetration-sharing sutta. The banyan-tree simile: defilements are self-generated, like the aerial roots of a banyan that grow from its own trunk.}

\section*{The Radiant Anthology: Non-Trainee Suttas}

And finally, the suttas that describe what lies beyond training --- the state of the arahant:

\begin{versebox}
Whose mind is like a rock ---\\
steady, unshaken ---\\
free of passion for what arouses passion,\\
unangered by what provokes anger ---\\
when one's mind is cultivated thus,\\
from where could suffering come?
\end{versebox}

\noindent Like a rock. The image is not of someone who doesn't feel, but of someone whose feeling doesn't move them off center.\footnote{Pali: \textit{``Yassa selūpamaṃ cittaṃ, ṭhitaṃ nānupakampati; virattaṃ rajanīyesu, kopaneyye na kuppati; yassevaṃ bhāvitaṃ cittaṃ, kuto naṃ dukkhamessatīti.''} (Ud 34). A non-trainee-sharing sutta.}

\begin{versebox}
Where water, earth,\\
fire and wind find no footing ---\\
there the stars do not shine,\\
the sun does not blaze,\\
the moon does not glow,\\
yet darkness is not found there.\\[6pt]
When the sage, the brahmin,\\
knows this for himself through silence ---\\
then from form and formless,\\
from happiness and pain, he is freed.
\end{versebox}

\noindent One of the most mysterious verses in the canon --- a place where the physical elements have no footing, where even light and darkness do not apply. Not a place, really --- a state beyond all categories.\footnote{Pali: \textit{``Yattha āpo ca pathavī, tejo vāyo na gādhati; na tattha sukkā jotanti, ādicco nappakāsati; na tattha candimā bhāti, tamo tattha na vijjati. Yadā ca attanāvedi, muni monena brāhmaṇo; atha rūpā arūpā ca, sukhadukkhā pamuccatīti.''} (Ud 10). A non-trainee-sharing sutta. This verse has generated extensive commentary. The Netti classifies it simply as describing the non-trainee's attainment.}

\begin{versebox}
Not by water is one pure ---\\
many people bathe here.\\
In whom there is truth and Dhamma ---\\
he is pure, he is a brahmin.
\end{versebox}

\begin{versebox}
When truths become manifest\\
to the ardent, meditating seeker ---\\
he stands scattering Māra's army\\
as the sun lights up the sky.
\end{versebox}

\noindent The same ``truths become manifest'' verse we met before --- but now with a different ending. Not ``doubts vanish'' (penetration), but ``scattering Māra's army like the sun'' (non-trainee). The work is complete. The battle is won.\footnote{Pali: ``Not by water'' = Ud 9 (\textit{na udakena sucī hoti}). ``Scattering Māra's army'' = Ud 1, third verse (\textit{vidhūpayaṃ tiṭṭhati mārasenaṃ, sūriyova obhāsayamantalikkhanti}). Both non-trainee-sharing suttas.}

\begin{versebox}
Not by slack effort,\\
not by small exertion,\\
is nirvana to be attained ---\\
the release from all suffering.\\[6pt]
This young monk,\\
this supreme person,\\
bears his final body,\\
having conquered Māra and his army.
\end{versebox}

\noindent And a prose passage that closes the non-trainee anthology:

``The Tathāgata, an arahant, fully awakened, is liberated through disenchantment with form, through its fading, through its cessation, through non-clinging. He is called `fully awakened.' A monk liberated through wisdom is liberated through disenchantment with form, through its fading, through its cessation, through non-clinging. He is called `liberated through wisdom.' What then is the distinction between them? The Tathāgata is the one who \textit{brought the unarisen path into being}, who \textit{produced the unproduced path}, who \textit{declared the undeclared path}. He is the knower of the path, the finder of the path, the one skilled in the path. And his disciples now dwell following that path. That is the distinction.''\footnote{Pali: \textit{``Tathāgato, bhikkhave, arahaṃ sammāsambuddho anuppannassa maggassa uppādetā, asañjātassa maggassa sañjanetā, anakkhātassa maggassa akkhātā, maggaññū maggavidū maggakovido, maggānugā ca bhikkhave etarahi sāvakā viharanti pacchāsamannāgatāti.''} (SN 22.58). The difference between a Buddha and an arahant: both are liberated, but the Buddha \textit{discovered} the path; the disciples \textit{follow} it. A non-trainee-sharing sutta. The Netti's classification system applies equally to both.}

\section*{The Combined Types}

Many suttas combine multiple types. Here is a defilement-and-habituation sutta:

\begin{versebox}
It rains right through a covered roof;\\
it does not rain through one laid open.\\
Therefore, open up the covered ---\\
and then it will not rain on you.
\end{versebox}

\noindent The ``covered'' mind (defilement) receives suffering; the ``opened'' mind (habituation toward practice) does not.\footnote{Pali: \textit{``Channamativassati, vivaṭaṃ nātivassatiti.''} (Ud 45). A combined defilement-and-habituation-sharing sutta: it names both the problem and the orientation toward solution (but not the specific path-analysis, which would make it penetration-sharing).}

The sixteen types exist not to create bureaucratic categories but to show that the Buddha's teaching is architecturally complete. Every possible angle is covered: the description of suffering, the accumulation of merit, the path of liberation, the state beyond training, and every combination of these.

\section*{The Eighteen Root Terms in the Suttas}

The Netti then maps the eighteen root terms onto specific examples. A \textbf{worldly} sutta: ``Evil deeds done do not curdle at once, like fresh milk; smoldering, they follow the fool like fire covered by ash.'' A \textbf{transcendent} sutta: one describing the destruction of the taints. A \textbf{being-based} sutta: one told in terms of persons. A \textbf{principle-based} sutta: one told in terms of mental qualities.\footnote{Pali: \textit{``Tattha katame aṭṭhārasa mūlapadā? Lokiyaṃ lokuttaraṃ lokiyañca lokuttarañca \ldots''} (§112). The worldly sutta quoted: \textit{``Na hi pāpaṃ kataṃ kammaṃ, sajjukhīraṃva muccati''} (Dhp 71). Each of the eighteen root terms is illustrated with canonical quotations, demonstrating that the classification system is not imposed from outside but discovered within the actual body of suttas.}

The point is completeness. Every discourse the Buddha ever spoke --- whether a single verse or a long prose sermon --- fits into at least one of the sixteen sutta-types and can be mapped to at least one of the eighteen root terms. And once mapped, it can be analyzed by any of the sixteen modes and interpreted according to any of the five methods.

The system has no outside. Nothing falls through the cracks.

\section*{The Four Grounds of Defilement}

The Netti closes with one last structural insight. All defilement-sharing suttas can be analyzed through four grounds of defilement: the ground of underlying tendency, the ground of active obsession, the ground of the fetter, and the ground of clinging.\footnote{Pali: \textit{``Saṃkilesabhāgiyaṃ suttaṃ catūhi kilesabhūmīhi niddisitabbaṃ --- anusayabhūmiyā pariyuṭṭhānabhūmiyā saṃyojanabhūmiyā upādānabhūmiyā.''} The four grounds of defilement (\textit{kilesabhūmi}): (1) underlying tendency (\textit{anusaya}) --- the dormant seed; (2) active obsession (\textit{pariyuṭṭhāna}) --- the seed sprouting; (3) fetter (\textit{saṃyojana}) --- the obsession binding; (4) clinging (\textit{upādāna}) --- the binding becoming fuel for existence.}

A being with underlying tendencies develops active obsession. The obsessed one becomes fettered. The fettered one clings. Dependent on clinging, existence. Dependent on existence, birth. Dependent on birth, aging, death, sorrow, lamentation, pain, grief, and despair. The whole mass of suffering.

That is the cascade --- and every defilement-sharing sutta in the entire canon is describing some part of it.

Habituation-sharing suttas should be analyzed through the three kinds of good conduct. Penetration-sharing suttas through the Four Noble Truths. Non-trainee-sharing suttas through the qualities of Buddhas, Paccekabuddhas, and noble disciples.\footnote{Pali: \textit{``Vāsanābhāgiyaṃ suttaṃ tīhi sucaritehi niddisitabbaṃ, nibbedhabhāgiyaṃ suttaṃ catūhi saccehi niddisitabbaṃ, asekkhabhāgiyaṃ suttaṃ tīhi dhammehi niddisitabbaṃ --- buddhadhammehi paccekabuddhadhammehi sāvakabhūmiyā. Jhāyivisaye niddisitabbanti.''}}

\ornament

The Netti concludes. The system that Mahākaccāyana built --- eighteen root terms, sixteen modes, five methods, sixteen types of sutta, and the four grounds of defilement --- is now complete.

You can take any passage from the Pali canon --- a verse from the Dhammapada, a discourse from the Majjhima Nikāya, a story from the Jātaka --- and run it through the Netti's system. The sixteen modes will show you its structure from sixteen different angles. The five methods will show you how it maps to the character-types and the liberation-doors. The sutta-classification will show you where it sits in the larger architecture. And the eighteen root terms will show you its deepest category.

The Netti is not, in the end, a book \textit{about} interpretation. It is a book that \textit{is} interpretation --- a live demonstration of how to read the words of the Buddha so that every sentence, every verse, every phrase yields its full meaning, in all its dimensions, with nothing left over and nothing missing.\footnote{Pali: \textit{``Nettippakaraṇaṃ niṭṭhitaṃ.''} ``The Nettippakaraṇa is completed.''}
